% !PRINTED 376
\weberpage{376}
%<<CONT
the constant factor in $\xi_{41}$ is suitably disposed of, also represent $f$ in the form:

\weq{7}{f = 4x_2\xi_3\,x_4\xi_{41} - v^2,}

in which $v$ is a quadratic form. From this and from the first representation in (5), however, there results the identity

\[
4x_2\xi_3(x_3\xi_2 - x_4\xi_{41}) = (u_1 - v)(u_1 + v),
\]

and from this we conclude that one of the two factors on the right must be divisible by $x_2\xi_3$ [\S.\,96, (3)]. If we assume that this is $u_1 - v$, which can be achieved by disposing of the sign of $v$, it follows that a constant factor $\lambda$ different from zero must exist, such that
% [sic] The source prints \lambda here; the following display and later prose use \lambda_1.

\[
\begin{aligned}
u_1 - v &= 2\lambda_1 x_2\xi_3,\\
u_1 + v &= \frac{2(x_3\xi_2 - x_4\xi_{41})}{\lambda_1},
\end{aligned}
\]

from which

\[
u_1 = \lambda_1 x_2\xi_3 + \frac{x_3\xi_2 - x_4\xi_{41}}{\lambda_1}
\]

follows. If we replace $u_1$ here by its expression from (6), we obtain

\[
x_4\xi_{41} = x_3\xi_2 - \lambda_1(-x_1\xi_1 + x_2\xi_2 + x_3\xi_3) + \lambda_1^2 x_2\xi_3.
\]

Of such equations we obtain, to begin with, three, if we permute the indices 1, 2, 3 cyclically and put in place of the unknown constant $\lambda_1$ three constants $\lambda_1, \lambda_2, \lambda_3$:

\weq{8}{\begin{aligned}
x_4\xi_{41} &= x_3\xi_2 - \lambda_1(-x_1\xi_1 + x_2\xi_2 + x_3\xi_3) + \lambda_1^2 x_2\xi_3\\
x_4\xi_{42} &= x_1\xi_3 - \lambda_2(x_1\xi_1 - x_2\xi_2 + x_3\xi_3) + \lambda_2^2 x_3\xi_1\\
x_4\xi_{43} &= x_2\xi_1 - \lambda_3(x_1\xi_1 + x_2\xi_2 - x_3\xi_3) + \lambda_3^2 x_1\xi_2.
\end{aligned}}

In order to determine the constants $\lambda_1, \lambda_2, \lambda_3$, we divide the last two of these equations by $\lambda_2$ and $\lambda_3$ and add. In this way follows the identical equation

\weq{9}{\begin{aligned}
x_4\left(\frac{\xi_{42}}{\lambda_2} + \frac{\xi_{43}}{\lambda_3}\right) &=\\
x_1\left(-2\xi_1 + \lambda_3\xi_2 + \frac{\xi_3}{\lambda_2}\right) &+ \xi_1\left(\lambda_2 x_3 + \frac{x_2}{\lambda_3}\right).
\end{aligned}}

Since now $x_4, x_1, \xi_1$ are azygetic, and consequently cannot pass through one point, the line

\[
\lambda_2 x_3 + \frac{x_2}{\lambda_3} = 0
\]
%CONT>>
